%----------------------------------------------------------
% 제12장. Model Retire와 Archiving
%----------------------------------------------------------

\chapter{Model Retire와 Archiving}
\label{ch:model_retirement}

AI 시스템의 생명주기는 배포와 운영으로 끝나지 않는다\cite{kreuzberger2023machine}. 모든 모델은 언젠가 Retire되며, 이 과정 역시 거버넌스의 중요한 영역이다. 체계적인 Retire 프로세스\cite{ashmore2021assuring}가 없으면 레거시 시스템이 방치되어 보안 위험이 되거나, 귀중한 지식이 소실되고, 규제 준수에 문제가 발생\cite{mokander2022conformity}한다. 특히 소유자가 불분명한 채 Resource만 소모하는 \textbf{'좀비 Model(Zombie Models)'}은 전사적 AI 거버넌스의 사각지대 중 하나이다 \cite{sculley2015hidden}.

본 장에서는 모델 Retire의 기준, 안전한 Retire 절차, Archiving 전략, 지식 보존\cite{sato2019continuous} 방법을 다룬다. 나아가 생성형 AI 시대의 LLM\cite{bommasani2021opportunities} 파인튜닝 모델 및 RAG\cite{lewis2020retrieval} 시스템 구성요소의 Retire 전략, 좀비 모델\cite{tamburri2020sustainable} 탐지 자동화, 그리고 EU AI Act가 요구하는 10년 문서 보존\cite{ebers2021european} 의무에 대응하는 실무 가이드를 제시한다.

%----------------------------------------------------------
\section{Model Retire 기준과 트리거}
\label{sec:12.1}

\subsection{Retirement 결정의 복잡성}
\label{sec:12.1.1}

Model Retirement 결정은 단순하지 않다. Technical적, Business적, 규제적 요인이 복합적으로 작용한다.

\begin{figure}[htbp]
\centering
\begin{tikzpicture}[
 node distance=1cm,
 factor/.style={rectangle, draw, rounded corners, minimum width=2.8cm, minimum height=1cm, align=center, font=\small},
 center/.style={circle, draw, fill=red!20, minimum size=2cm, align=center, font=\small\bfseries},
]

% 중앙
\node[center] (decision) at (0, 0) {Retire\\결정};

% 4대 요인
\node[factor, fill=blue!20] (tech) at (0, 3) {\textbf{Technical적 요인}\\성능 Degradation, Technical 노후화};
\node[factor, fill=green!20] (biz) at (3.5, 0) {\textbf{Business 요인}\\가치 상실, 전략 변화};
\node[factor, fill=orange!20] (reg) at (0, -3) {\textbf{규제 요인}\\Regulatory Change, 컴플라이언스};
\node[factor, fill=purple!20] (ops) at (-3.5, 0) {\textbf{Operational 요인}\\비용, 유지보수 부담};

% 연결
\draw[->, thick] (tech) -- (decision);
\draw[->, thick] (biz) -- (decision);
\draw[->, thick] (reg) -- (decision);
\draw[->, thick] (ops) -- (decision);

\end{tikzpicture}
\caption{Model Retirement 결정의 4대 요인}
\label{fig:retirement_factors}
\end{figure}

특히 정기적인 \textbf{'Model 건강검진'}을 통해 6개월 이상 사용량이 없거나 ROI가 급감한 '좀비 Model'을 탐지하는 프로세스가 12.1.3절의 프레임워크에 포함되어야 한다.\cite{shankar2022operationalizing}

\subsection{전사 리소스 최적화 KPI와 ROI}
\label{sec:12.1.3_kpi}

좀비 모델의 방치는 단순히 기술적 부채를 넘어 경영 효율성에 직접적인 타격을 준다. 이를 관리하기 위해 다음과 같은 \textbf{리소스 최적화 KPI}를 전사 거버넌스 지표에 관리할 것을 권고한다.

\begin{itemize}
 \item \textbf{인프라 가용성 회수율(인프라 Reclamation Rate)}: Retire를 통해 회수된 GPU/CPU 자원 및 스토리지의 비율. 이는 클라우드 비용 절감으로 직결된다.
 \item \textbf{운영 복잡도 감소 지수(Operational Complexity Reduction)}: 관리 대상 모델 수의 최적화를 통해 운영 인력의 1인당 관리 모델 수를 적정 수준으로 유지.
 \item \textbf{AI 자산 회전율(AI Asset Turnover)}: 개발된 모델 중 실제 비즈니스 가치를 창출하며 운영 중인 모델의 비중.
\end{itemize}

이러한 지표들은 경영진에게 AI 거버넌스가 단순한 '규제 대응'이 아니라 \textbf{'비용 최적화와 이익 극대화'}를 위한 전략적 도구임을 입증하는 핵심 근거가 된다.

\subsection{Retire 트리거 분류}
\label{sec:12.1.2}

\begin{table}[htbp]
\centering
\caption{Model Retire 트리거 분류}
\label{tab:retirement_triggers}
\begin{tabularx}{\textwidth}{p{2.5cm}p{3.5cm}X}
\toprule
\textbf{카테고리} & \textbf{트리거} & \textbf{설명 및 결정 기준} \\
\midrule
\multirow{4}{*}{\textbf{Technical적}} &
성능 Degradation (드리프트\cite{paleyes2022challenges}) &
모니터링 KPI가 임계값 이하로 지속 하락;
재학습으로 회복 불가 \\
\cmidrule{2-3}
& Model Obsolescence &
더 나은 알고리즘/아키텍처 등장;
경쟁력 상실 \\
\cmidrule{2-3}
& Tech Stack EOL &
의존 프레임워크/라이브러리 Support 종료 \\
\cmidrule{2-3}
& Security Vulnerability &
해결 불가능한 보안 이슈 발견 \\
\midrule
\multirow{3}{*}{\textbf{Business}} &
Business Value Loss &
ROI 미달; 사용량 급감;
대상 업무 폐지 \\
\cmidrule{2-3}
& 전략 Change &
사업 방향 Migration;
해당 영역 철수 \\
\cmidrule{2-3}
& Replacement Ready &
신규 Model/시스템으로 완전 대체 \\
\midrule
\multirow{3}{*}{\textbf{규제}} &
Regulatory Change &
관련 규제 변경으로 준수 불가;
금지 대상 지정 \\
\cmidrule{2-3}
& 컴플라이언스 Failure &
감사 지적; 규제 위반 확인 \\
\cmidrule{2-3}
& Ethical Issue &
심각한 편향, 차별 문제 확인;
해결 불가 \\
\midrule
\multirow{3}{*}{\textbf{Operational}} &
Maintenance Cost &
운영 비용이 가치 초과;
전문 인력 부재 \\
\cmidrule{2-3}
& 인프라 Constraint &
필요 인프라 Support 종료;
마이그레이션 불가 \\
\cmidrule{2-3}
& Integration Complexity &
시스템 간 Integration 문제 누적;
Technical 부채 과다\cite{sculley2015hidden} \\
\bottomrule
\end{tabularx}
\end{table}


\subsection{Retirement Decision 프레임워크}
\label{sec:12.1.3}

\begin{lstlisting}[language=Python, caption={Model Retirement Decision Framework}, label={lst:retirement_decision}]
"""
Model Retirement Decision Framework - Trigger Detection and Decision Support
"""

from dataclasses import dataclass, field
from typing import Dict, List, Optional, Tuple
from datetime import date, datetime, timedelta
from enum import Enum
import numpy as np

class RetirementTrigger(Enum):
 """Retirement Trigger Type"""
 # Technical
 PERFORMANCE_DEGRADATION = "Performance Degradation"
 MODEL_OBSOLESCENCE = "Model Obsolescence"
 TECH_STACK_EOL = "Tech Stack EOL"
 SECURITY_VULNERABILITY = "Security Vulnerability"

 # Business
 BUSINESS_VALUE_LOSS = "Business Value Loss"
 STRATEGY_CHANGE = "Strategy Change"
 REPLACEMENT_READY = "Replacement Ready"

 # Regulation
 REGULATORY_CHANGE = "Regulatory Change"
 COMPLIANCE_FAILURE = "Compliance Failure"
 ETHICAL_ISSUE = "Ethical Issue"

 # Operational
 MAINTENANCE_COST = "Maintenance Cost"
 INFRASTRUCTURE_CONSTRAINT = "Infrastructure Constraint"
 INTEGRATION_COMPLEXITY = "Integration Complexity"

class RetirementUrgency(Enum):
 """Retirement Urgency"""
 IMMEDIATE = "Immediate Retirement" # 24 URGENT = "Urgent Retirement" # 1 PLANNED = "Planned Retirement" # 1-3months
 SCHEDULED = "Scheduled Retirement" # 3months

class RetirementDecision(Enum):
 """Retirement Decision"""
 RETIRE = "Retire"
 REMEDIATE = "Remediate and Maintain"
 MONITOR = "Continue Monitoring"
 DEFER = "Defer Decision"

@dataclass
class ModelHealthMetrics:
 """Model Health Metrics"""
 model_id: str
 assessment_date: date

 # Performance Metric
 current_accuracy: float = 0.0
 baseline_accuracy: float = 0.0
 accuracy_trend: str = "stable" # improving, stable, declining
 drift_score: float = 0.0 # 0-1, 1 # Business Metric
 daily_predictions: int = 0
 usage_trend: str = "stable"
 estimated_monthly_value: float = 0.0
 estimated_monthly_cost: float = 0.0

 # Operational Metric
 incidents_last_90days: int = 0
 avg_latency_ms: float = 0.0
 uptime_percent: float = 99.9
 maintenance_hours_monthly: float = 0.0

 # Compliance Metric
 last_audit_date: Optional[date] = None
 audit_issues: int = 0
 documentation_complete: bool = True

 # Technical Metric
 framework_version: str = ""
 framework_eol_date: Optional[date] = None
 security_vulnerabilities: int = 0

@dataclass
class RetirementAssessment:
 """Retire Assessment """
 model_id: str
 assessment_date: date
 assessor: str

 #
 triggered_by: List[RetirementTrigger] = field(default_factory=list)

 #
 technical_score: float = 0.0 # 0-100, business_score: float = 0.0
 compliance_score: float = 0.0
 operational_score: float = 0.0
 overall_score: float = 0.0

 #
 recommended_decision: RetirementDecision = RetirementDecision.MONITOR
 urgency: RetirementUrgency = RetirementUrgency.SCHEDULED
 rationale: str = ""

 # Impact
 affected_systems: List[str] = field(default_factory=list)
 affected_users: int = 0


class ModelRetirementFramework:
 """Model Retirement Decision Framework"""

 # Setup
 THRESHOLDS = {
 'accuracy_drop_critical': 0.15, # Critical: 15% drop from baseline
 'accuracy_drop_warning': 0.10, # Warning: 10% drop
 'drift_critical': 0.3, # Critical: Drift score >= 0.3
 'drift_warning': 0.2,
 'usage_drop_critical': 0.7, # Critical: 70% decrease in usage
 'roi_minimum': 1.0, # Re-evaluate if ROI < 1.0
 'incident_critical': 5, # Critical: >= 5 incidents in 90 days
 'vulnerability_critical': 1, # Security Vulnerability 1 }

 def __init__(self):
 pass

 def assess_retirement(
 self,
 metrics: ModelHealthMetrics,
 assessor: str,
 ) -> RetirementAssessment:
 """Retire Assessment """

 assessment = RetirementAssessment(
 model_id=metrics.model_id,
 assessment_date=date.today(),
 assessor=assessor,
 )

 triggers = []

 # 1. Technical Assessment
 tech_score = 100

 # Performance Degradation
 if metrics.baseline_accuracy > 0:
 accuracy_drop = (metrics.baseline_accuracy - metrics.current_accuracy) / metrics.baseline_accuracy
 if accuracy_drop >= self.THRESHOLDS['accuracy_drop_critical']:
 triggers.append(RetirementTrigger.PERFORMANCE_DEGRADATION)
 tech_score -= 40
 elif accuracy_drop >= self.THRESHOLDS['accuracy_drop_warning']:
 tech_score -= 20

 # Check Drift
 if metrics.drift_score >= self.THRESHOLDS['drift_critical']:
 triggers.append(RetirementTrigger.PERFORMANCE_DEGRADATION)
 tech_score -= 30
 elif metrics.drift_score >= self.THRESHOLDS['drift_warning']:
 tech_score -= 15

 # Check Framework EOL
 if metrics.framework_eol_date:
 days_to_eol = (metrics.framework_eol_date - date.today()).days
 if days_to_eol <= 0:
 triggers.append(RetirementTrigger.TECH_STACK_EOL)
 tech_score -= 30
 elif days_to_eol <= 90:
 tech_score -= 15

 # Security Vulnerability
 if metrics.security_vulnerabilities >= self.THRESHOLDS['vulnerability_critical']:
 triggers.append(RetirementTrigger.SECURITY_VULNERABILITY)
 tech_score -= 40

 assessment.technical_score = max(0, tech_score)

 # 2. Business Assessment
 biz_score = 100

 # Check ROI
 if metrics.estimated_monthly_cost > 0:
 roi = metrics.estimated_monthly_value / metrics.estimated_monthly_cost
 if roi < self.THRESHOLDS['roi_minimum']:
 triggers.append(RetirementTrigger.BUSINESS_VALUE_LOSS)
 biz_score -= 40

 # Check Usage
 if metrics.usage_trend == "declining":
 biz_score -= 20

 assessment.business_score = max(0, biz_score)

 # 3. Compliance Assessment
 compliance_score = 100

 if metrics.audit_issues > 0:
 triggers.append(RetirementTrigger.COMPLIANCE_FAILURE)
 compliance_score -= metrics.audit_issues * 20

 if not metrics.documentation_complete:
 compliance_score -= 20

 assessment.compliance_score = max(0, compliance_score)

 # 4. Operational Assessment
 ops_score = 100

 if metrics.incidents_last_90days >= self.THRESHOLDS['incident_critical']:
 ops_score -= 30

 if metrics.uptime_percent < 99.0:
 ops_score -= 20

 assessment.operational_score = max(0, ops_score)

 # 5. Overall Score and Decision
 assessment.overall_score = (
 assessment.technical_score * 0.3 +
 assessment.business_score * 0.3 +
 assessment.compliance_score * 0.25 +
 assessment.operational_score * 0.15
 )

 assessment.triggered_by = triggers

 # Logic
 assessment.recommended_decision, assessment.urgency = self._determine_decision(
 assessment, triggers
 )

 assessment.rationale = self._generate_rationale(assessment, triggers)

 return assessment

 def _determine_decision(
 self,
 assessment: RetirementAssessment,
 triggers: List[RetirementTrigger],
 ) -> Tuple[RetirementDecision, RetirementUrgency]:
 """Retirement Decision """

 # Immediate Retirement
 immediate_triggers = [
 RetirementTrigger.SECURITY_VULNERABILITY,
 RetirementTrigger.COMPLIANCE_FAILURE,
 RetirementTrigger.ETHICAL_ISSUE,
 ]

 if any(t in triggers for t in immediate_triggers):
 return RetirementDecision.RETIRE, RetirementUrgency.IMMEDIATE

 # Urgent Retirement
 if RetirementTrigger.TECH_STACK_EOL in triggers:
 return RetirementDecision.RETIRE, RetirementUrgency.URGENT

 # Score-based Decision
 if assessment.overall_score < 40:
 return RetirementDecision.RETIRE, RetirementUrgency.PLANNED
 elif assessment.overall_score < 60:
 return RetirementDecision.REMEDIATE, RetirementUrgency.PLANNED
 elif assessment.overall_score < 80:
 return RetirementDecision.MONITOR, RetirementUrgency.SCHEDULED
 else:
 return RetirementDecision.DEFER, RetirementUrgency.SCHEDULED

 def _generate_rationale(
 self,
 assessment: RetirementAssessment,
 triggers: List[RetirementTrigger],
 ) -> str:
 """Generate Decision Rationale"""

 rationale_parts = []

 rationale_parts.append(f"Overall Health Score: {assessment.overall_score:.1f}/100")

 if triggers:
 trigger_names = [t.value for t in triggers]
 rationale_parts.append(f"Active Triggers: {', '.join(trigger_names)}")

 # Weakest Area
 scores = {
 'Technical': assessment.technical_score,
 'Business': assessment.business_score,
 'Compliance': assessment.compliance_score,
 'Operational': assessment.operational_score,
 }
 weakest = min(scores, key=scores.get)
 rationale_parts.append(f"Weakest Area: {weakest} ({scores[weakest]:.1f} )")

 return "; ".join(rationale_parts)

 def generate_retirement_report(self, assessment: RetirementAssessment) -> str:
 """Retire Assessment Report Generation - ( Guide Retention)"""
 # ( : Code days Report Generation Logic)
 return "# Model Retire Assessment Report ( )"
\end{lstlisting}


\subsection{좀비 모델 탐지 자동화}
\label{sec:12.1.4}

좀비 모델(Zombie Model)은 공식적으로 Retire되지 않았으나 실질적으로 가치를 창출하지 못하는 모델을 의미한다. 2025년 기준 대규모 AI 도입 기업의 약 30\%가 소유자 불분명, 사용량 제로, 또는 성능 미달 상태의 좀비 모델을 10개 이상 운영 중인 것으로 추정된다\cite{sculley2015hidden}. 이러한 모델들은 다음과 같은 비용을 발생시킨다.

\begin{itemize}
 \item \textbf{인프라 비용}: GPU/CPU 자원 점유, 클라우드 인스턴스 유지 비용
 \item \textbf{보안 리스크}: 패치되지 않는 레거시 프레임워크의 취약점 노출
 \item \textbf{거버넌스 사각지대}: 모니터링, 감사, 편향 점검이 이루어지지 않는 모델
 \item \textbf{운영 혼란}: 신규 인력이 해당 모델의 역할과 소유자를 파악하지 못함
\end{itemize}

다음 코드는 모델 레지스트리와 인프라 모니터링 데이터를 기반으로 좀비 모델을 자동 탐지하는 시스템을 구현한 것이다.

\begin{lstlisting}[language=Python, caption={좀비 모델 탐지 자동화 시스템}, label={lst:zombie_detection}]
"""
Zombie Model Detection System
- 사용량, 성능, 비용, 소유자 메트릭을 기반으로 좀비 모델 자동 탐지
- 정기 스캔(주간/월간)으로 레지스트리 내 전체 모델을 점검
"""

from dataclasses import dataclass, field
from datetime import date, timedelta
from typing import List, Dict, Optional
from enum import Enum
import logging

logger = logging.getLogger(__name__)

class ZombieRiskLevel(Enum):
    """좀비 위험 등급"""
    CRITICAL = "Critical"    # 즉시 Retire 권고
    HIGH = "High"            # 30일 내 소유자 확인 필요
    MEDIUM = "Medium"        # 모니터링 강화
    LOW = "Low"              # 정상

@dataclass
class ModelRegistryEntry:
    """모델 레지스트리 항목"""
    model_id: str
    model_name: str
    owner: Optional[str]
    team: Optional[str]
    deployed_date: date
    last_prediction_date: Optional[date]
    framework: str
    framework_version: str
    endpoint_url: Optional[str] = None
    tags: List[str] = field(default_factory=list)

@dataclass
class ModelUsageMetrics:
    """모델 사용량 메트릭"""
    model_id: str
    daily_request_count_avg_30d: float = 0.0
    daily_request_count_avg_90d: float = 0.0
    unique_consumers_30d: int = 0
    last_request_timestamp: Optional[date] = None

@dataclass
class ModelCostMetrics:
    """모델 비용 메트릭"""
    model_id: str
    monthly_compute_cost: float = 0.0
    monthly_storage_cost: float = 0.0
    monthly_network_cost: float = 0.0
    estimated_monthly_value: float = 0.0

    @property
    def total_monthly_cost(self) -> float:
        return (self.monthly_compute_cost
                + self.monthly_storage_cost
                + self.monthly_network_cost)

    @property
    def roi(self) -> float:
        if self.total_monthly_cost == 0:
            return 0.0
        return self.estimated_monthly_value / self.total_monthly_cost

@dataclass
class ZombieDetectionResult:
    """좀비 탐지 결과"""
    model_id: str
    risk_level: ZombieRiskLevel
    reasons: List[str] = field(default_factory=list)
    days_since_last_use: Optional[int] = None
    monthly_waste_cost: float = 0.0
    recommended_action: str = ""

class ZombieModelDetector:
    """
    좀비 모델 탐지기

    탐지 기준:
    1. 사용량 기준: 90일간 평균 요청 수 < threshold
    2. 소유자 기준: owner 필드가 비어 있거나 퇴사자
    3. 비용 기준: ROI < 1.0 (비용 > 가치)
    4. 보안 기준: EOL된 프레임워크 사용
    5. 기간 기준: 마지막 사용일로부터 N일 이상 경과
    """

    DEFAULT_CONFIG = {
        'idle_days_critical': 180,      # 180일 미사용 -> Critical
        'idle_days_warning': 90,        # 90일 미사용 -> High
        'min_daily_requests': 1.0,      # 일 평균 1건 미만 -> 의심
        'min_unique_consumers': 1,      # 소비자 0 -> 의심
        'roi_threshold': 1.0,           # ROI 1.0 미만 -> 비용 초과
        'eol_frameworks': [             # EOL된 프레임워크 목록
            'tensorflow-1.x',
            'pytorch-1.x',
            'scikit-learn-0.x',
        ],
    }

    def __init__(self, config: Optional[Dict] = None):
        self.config = config or self.DEFAULT_CONFIG
        self.departed_employees: List[str] = []  # HR 연동

    def scan_registry(
        self,
        registry: List[ModelRegistryEntry],
        usage_data: Dict[str, ModelUsageMetrics],
        cost_data: Dict[str, ModelCostMetrics],
    ) -> List[ZombieDetectionResult]:
        """전체 레지스트리 스캔"""
        results = []

        for entry in registry:
            result = self._evaluate_model(
                entry,
                usage_data.get(entry.model_id),
                cost_data.get(entry.model_id),
            )
            if result.risk_level != ZombieRiskLevel.LOW:
                results.append(result)
                logger.warning(
                    f"Zombie detected: {entry.model_id} "
                    f"[{result.risk_level.value}] "
                    f"- {', '.join(result.reasons)}"
                )

        # 위험도 순 정렬
        results.sort(
            key=lambda r: list(ZombieRiskLevel).index(r.risk_level)
        )
        return results

    def _evaluate_model(
        self,
        entry: ModelRegistryEntry,
        usage: Optional[ModelUsageMetrics],
        cost: Optional[ModelCostMetrics],
    ) -> ZombieDetectionResult:
        """개별 모델 좀비 여부 평가"""
        reasons = []
        risk_scores = []
        monthly_waste = 0.0
        days_idle = None

        # 1. 사용량 점검
        if usage:
            if usage.last_request_timestamp:
                days_idle = (date.today()
                             - usage.last_request_timestamp).days
                if days_idle >= self.config['idle_days_critical']:
                    reasons.append(
                        f"{days_idle}일간 미사용 (Critical)")
                    risk_scores.append(3)
                elif days_idle >= self.config['idle_days_warning']:
                    reasons.append(
                        f"{days_idle}일간 미사용 (Warning)")
                    risk_scores.append(2)

            if (usage.daily_request_count_avg_90d
                    < self.config['min_daily_requests']):
                reasons.append("일 평균 요청 수 1건 미만")
                risk_scores.append(2)

            if (usage.unique_consumers_30d
                    < self.config['min_unique_consumers']):
                reasons.append("활성 소비자 없음")
                risk_scores.append(2)
        else:
            reasons.append("사용량 데이터 없음")
            risk_scores.append(2)

        # 2. 소유자 점검
        if not entry.owner:
            reasons.append("소유자 미지정")
            risk_scores.append(2)
        elif entry.owner in self.departed_employees:
            reasons.append(
                f"소유자 퇴사: {entry.owner}")
            risk_scores.append(2)

        # 3. 비용 점검
        if cost:
            if cost.roi < self.config['roi_threshold']:
                reasons.append(
                    f"ROI {cost.roi:.2f} (임계값 미만)")
                risk_scores.append(2)
                monthly_waste = (cost.total_monthly_cost
                                 - cost.estimated_monthly_value)

        # 4. 프레임워크 EOL 점검
        fw_key = (f"{entry.framework}-"
                  f"{entry.framework_version.split('.')[0]}.x")
        if fw_key in self.config['eol_frameworks']:
            reasons.append(
                f"EOL 프레임워크: {entry.framework} "
                f"{entry.framework_version}")
            risk_scores.append(2)

        # 종합 위험도 결정
        if not risk_scores:
            risk_level = ZombieRiskLevel.LOW
        elif max(risk_scores) >= 3 or len(risk_scores) >= 3:
            risk_level = ZombieRiskLevel.CRITICAL
        elif max(risk_scores) >= 2:
            risk_level = ZombieRiskLevel.HIGH
        else:
            risk_level = ZombieRiskLevel.MEDIUM

        # 권고 액션 결정
        action_map = {
            ZombieRiskLevel.CRITICAL:
                "즉시 Retire 프로세스 개시",
            ZombieRiskLevel.HIGH:
                "30일 내 소유자 확인 및 Retire 검토",
            ZombieRiskLevel.MEDIUM:
                "모니터링 강화 및 분기 내 재평가",
            ZombieRiskLevel.LOW: "정상 운영",
        }

        return ZombieDetectionResult(
            model_id=entry.model_id,
            risk_level=risk_level,
            reasons=reasons,
            days_since_last_use=days_idle,
            monthly_waste_cost=max(0, monthly_waste),
            recommended_action=action_map[risk_level],
        )

    def generate_summary_report(
        self,
        results: List[ZombieDetectionResult],
    ) -> Dict:
        """좀비 탐지 요약 리포트 생성"""
        summary = {
            'scan_date': date.today().isoformat(),
            'total_scanned': 0,
            'zombies_found': len(results),
            'by_risk_level': {},
            'total_monthly_waste': sum(
                r.monthly_waste_cost for r in results),
            'top_waste_models': [],
        }

        for level in ZombieRiskLevel:
            count = sum(
                1 for r in results
                if r.risk_level == level)
            summary['by_risk_level'][level.value] = count

        # 비용 낭비 상위 5개 모델
        sorted_by_waste = sorted(
            results,
            key=lambda r: r.monthly_waste_cost,
            reverse=True,
        )[:5]
        summary['top_waste_models'] = [
            {
                'model_id': r.model_id,
                'monthly_waste': r.monthly_waste_cost,
                'reasons': r.reasons,
            }
            for r in sorted_by_waste
        ]

        return summary
\end{lstlisting}

\paragraph{좀비 탐지 자동화 운영 가이드}
좀비 모델 탐지 시스템은 다음과 같은 주기로 운영한다.

\begin{itemize}
 \item \textbf{주간 스캔}: 사용량 기반 경량 스캔으로 Critical 등급 모델 즉시 발견
 \item \textbf{월간 스캔}: 비용, 소유자, 프레임워크 EOL을 포함한 종합 스캔
 \item \textbf{분기 리포트}: 경영진에게 좀비 모델 현황 및 비용 절감 효과 보고
 \item \textbf{자동 알림}: Critical 등급 탐지 시 Slack/Teams 채널에 즉시 알림 전송
\end{itemize}


%----------------------------------------------------------
\section{안전한 Retire 프로세스}
\label{sec:12.2}

\subsection{Retire 프로세스 개요}
\label{sec:12.2.1}

Model Retire는 체계적인 프로세스를 따라야 한다. 성급한 Retire는 서비스 장애, 데이터 손실, 규제 위반을 초래할 수 있다.

\begin{figure}[htbp]
\centering
\begin{tikzpicture}[
 node distance=0.8cm,
 phase/.style={rectangle, draw, rounded corners, minimum width=2.5cm, minimum height=1cm, align=center, font=\small},
 arrow/.style={->, thick},
]

% 단계
\node[phase, fill=blue!20] (assess) at (0, 0) {1. Retire\\평가};
\node[phase, fill=blue!20] (approve) at (3, 0) {2. Approved\\획득};
\node[phase, fill=green!20] (notify) at (6, 0) {3. 이해관계자\\Notification};
\node[phase, fill=green!20] (migrate) at (9, 0) {4. 마이그레이션\\/Migration};
\node[phase, fill=orange!20] (archive) at (12, 0) {5. Archiving};

\node[phase, fill=orange!20] (shutdown) at (3, -2) {6. 서비스\\Shutdown};
\node[phase, fill=red!20] (cleanup) at (6, -2) {7. Resource\\회수};
\node[phase, fill=red!20] (document) at (9, -2) {8. 최종\\Documented};
\node[phase, fill=gray!20] (close) at (12, -2) {9. Retire\\Completed};

% 화살표
\draw[arrow] (assess) -- (approve);
\draw[arrow] (approve) -- (notify);
\draw[arrow] (notify) -- (migrate);
\draw[arrow] (migrate) -- (archive);
\draw[arrow] (archive) -- (shutdown);
\draw[arrow] (shutdown) -- (cleanup);
\draw[arrow] (cleanup) -- (document);
\draw[arrow] (document) -- (close);

\end{tikzpicture}
\caption{Model Retire 프로세스}
\label{fig:retirement_process}
\end{figure}


\begin{table}[htbp]
\centering
\caption{Retire 프로세스 단계별 상세}
\label{tab:retirement_process_details}
\begin{tabularx}{\textwidth}{p{0.5cm}p{2cm}p{2.5cm}X}
\toprule
\textbf{} & \textbf{단계} & \textbf{담당} & \textbf{주요 활동 및 산출물} \\
\midrule
1 &
Retire 평가 &
AI 거버넌스팀 &
트리거 확인, 건강도 평가, 영향 분석 -> Retire 평가 Report \\
\midrule
2 &
Approved 획득 &
AI 위원회/CDO &
평가 검토, Retirement Decision Approved -> Retire Approved서 \\
\midrule
3 &
이해관계자 Notification &
프로젝트 매니저 &
사용 부서, Integration 시스템 담당자 Notification -> Notification 기록, 협의 결과 \\
\midrule
4 &
마이그레이션/Migration &
개발팀 &
대체 시스템으로 트래픽 Migration, 데이터 마이그레이션 -> Migration Completed 보고 \\
\midrule
5 &
Archiving &
MLOps팀 &
모델, 데이터, Archive Documentation\cite{zaharia2018accelerating} -> 아카이브\cite{huyen2022designing} 목록 \\
\midrule
6 &
서비스 Shutdown &
인프라팀 &
API 비활성화, 엔드포인트 차단 -> Shutdown 확인 \\
\midrule
7 &
Resource 회수 &
인프라팀 &
컴퓨팅 Resource, 스토리지 회수 -> Resource 회수 목록 \\
\midrule
8 &
최종 Documented &
AI 거버넌스팀 &
Retire 경위, 학습 사항 Documented -> Retire Completed Report \\
\midrule
9 &
Retire Completed &
AI 거버넌스팀 &
레지스트리 상태 갱신, 종결 -> 상태: Retire됨 \\
\bottomrule
\end{tabularx}
\end{table}


\subsection{긴급도별 Retire 절차}
\label{sec:12.2.2}

Retirement Urgency에 따라 절차의 속도와 범위가 달라진다. 특히 보안 Vulnerability Attack 성공 등 치명적 리스크 발생 시에는 이해관계자 사후 Notification를 허용하는 'Immediate Retirement' 트랙을 운영한다.

\begin{table}[htbp]
\centering
\caption{긴급도별 Retire 절차}
\label{tab:urgency_procedures}
\begin{tabularx}{\textwidth}{p{2cm}p{2cm}X}
\toprule
\textbf{긴급도} & \textbf{목표 기간} & \textbf{절차 특이사항} \\
\midrule
Immediate Retirement &
24시간 이내 &
 Security/컴플라이언스 긴급 Response\\
& & 즉시 API 차단 가능\\
& & 사후 이해관계자 Notification 허용\\
& & 최소 Archiving 후 Retire \\
\midrule
Urgent Retirement &
1주 이내 &
 단축 Approved 절차\\
& & 주요 이해관계자만 사전 Notification\\
& & 병행 작업 (Archiving + Migration)\\
\midrule
Planned Retirement &
1-3개월 &
 표준 절차 컴플라이언스\\
& & 충분한 Migration 기간\\
& & 완전한 Archiving 및 Documented \\
\bottomrule
\end{tabularx}
\end{table}


\subsection{Retire 관리 System}
\label{sec:12.2.3}

\begin{lstlisting}[language=Python, caption={Model Retirement Management System Implementation}, label={lst:retirement_management}]
class RetirementPhase(Enum):
 INITIATED = "Initiated"; ASSESSMENT = "Assessment"; APPROVED = "Approved"
 NOTIFICATION = "Notification"; MIGRATION = "Migration"; ARCHIVING = "Archiving"
 SHUTDOWN = "Shutdown"; CLEANUP = "Cleanup"; DOCUMENTED = "Documented"
 COMPLETED = "Completed"

@dataclass
class RetirementProject:
 project_id: str; model_id: str
 current_phase: RetirementPhase = RetirementPhase.INITIATED
 urgency: RetirementUrgency = RetirementUrgency.PLANNED
 tasks: List[RetirementTask] = field(default_factory=list)
 archive_location: str = ""

class ModelRetirementManager:
 """Model Retire Workflow Management """
 PHASE_TASKS = {
 RetirementPhase.ARCHIVING: ["Archive Model Artifacts", "Archive Training Data References", "Archive Documentation"],
 RetirementPhase.SHUTDOWN: ["Deactivate API Endpoints", "Disable Monitoring Alerts"]
 }

 def advance_phase(self, project_id: str) -> RetirementProject:
 # Transition to next phase after checking task completion
 pass
\end{lstlisting}


%----------------------------------------------------------
\section{Archiving 전략}
\label{sec:12.3}

\subsection{아카이브 대상 및 Scope}
\label{sec:12.3.1}

\begin{table}[htbp]
\centering
\caption{아카이브 대상 분류 및 Retention Period}
\label{tab:archive_scope}
\small
\begin{tabularx}{\textwidth}{p{2cm}p{2.5cm}X p{2.5cm}}
\toprule
\textbf{카테고리} & \textbf{Artifact} & \textbf{보존 근거} & \textbf{Retention Period} \\
\midrule
\multirow{2}{*}{\textbf{Model}} &
Model Weights/코드 &
재현성, 감사 Response &
5years/Permanent \\
\cmidrule{2-4}
& 하이퍼파라미터 &
지식 보존 &
5years \\
\midrule
\multirow{2}{*}{\textbf{Data}} &
데이터 참조/피처 &
리니지 추적 (9장 연계) &
5years \\
\cmidrule{2-4}
& 평가 데이터셋 &
벤치마크 비교 &
3years \\
\midrule
\multirow{2}{*}{\textbf{문서}} &
Model Card/평가서 &
규제 준수 (EU AI Act) &
5years \\
\cmidrule{2-4}
& Approved 기록 &
감사 증빙 &
5years \\
\bottomrule
\end{tabularx}
\end{table}

\subsection{클라우드 vs 온프레미스 Archiving 전략}
\label{sec:12.3.1b}

아카이브 인프라의 선택은 비용, 규제 요건, 접근 빈도에 따라 결정된다. 2025년 현재 대부분의 기업은 하이브리드 전략을 채택하며, 핵심은 \textbf{계층형 스토리지(Tiered Storage)} 아키텍처의 설계이다.

\begin{table}[htbp]
\centering
\caption{클라우드 vs 온프레미스 Archiving 비교}
\label{tab:cloud_vs_onprem_archive}
\begin{tabularx}{\textwidth}{p{2.5cm}XX}
\toprule
\textbf{항목} & \textbf{클라우드 Archiving} & \textbf{온프레미스 Archiving} \\
\midrule
\textbf{비용 구조} &
종량제(Pay-as-you-go); 초기 투자 최소화; 장기 보관 시 Glacier/Archive Tier 활용 &
초기 하드웨어 투자 필요; 장기적으로 대용량 시 유리; 전력 및 유지보수 비용 발생 \\
\midrule
\textbf{확장성} &
사실상 무제한; 자동 확장; 멀티 리전 복제 용이 &
물리적 한계; 확장 시 조달 리드타임; 용량 계획 필수 \\
\midrule
\textbf{규제 준수} &
데이터 주권 이슈 발생 가능; 리전 선택으로 대응; GDPR/개인정보보호법 검토 필요 &
데이터 주권 확보 용이; 물리적 접근 통제; 규제 기관 현장 감사 대응 유리 \\
\midrule
\textbf{접근 속도} &
Hot: 즉시 / Warm: 수 분 / Cold: 수 시간; Archive Tier: 12--48시간 &
Hot: 즉시 / Cold: 테이프 복원 시 수 시간; 물리적 위치에 따라 편차 \\
\midrule
\textbf{보안} &
CSP 관리형 암호화; IAM 통합; 감사 로그 자동화 &
자체 암호화 키 관리; 물리적 보안; 네트워크 격리 \\
\midrule
\textbf{권장 시나리오} &
다수 모델의 장기 보존; 글로벌 조직; 가변적 접근 패턴 &
고위험 AI(의료, 금융); 데이터 주권 필수; 대용량 학습 데이터 보존 \\
\bottomrule
\end{tabularx}
\end{table}

\paragraph{계층형 스토리지 아키텍처}
AI 모델 아카이브를 위한 계층형 스토리지는 접근 빈도와 규제 요건에 따라 다음 3개 계층으로 구성한다.

\begin{figure}[htbp]
\centering
\begin{tikzpicture}[
 tier/.style={rectangle, draw, rounded corners, minimum width=10cm, minimum height=1.2cm, align=center, font=\small},
]

\node[tier, fill=red!20] (hot) at (0, 3) {\textbf{Hot Tier} (NVMe/SSD) -- 최근 Retire된 모델 (0--90일)\\즉시 접근, 감사 대응, 롤백 대비};
\node[tier, fill=yellow!20] (warm) at (0, 1.2) {\textbf{Warm Tier} (S3 Standard/Azure Blob Hot) -- 중기 보존 (90일--2년)\\분 단위 접근, 벤치마크 비교, 지식 참조};
\node[tier, fill=blue!15] (cold) at (0, -0.6) {\textbf{Cold Tier} (S3 Glacier/Azure Archive) -- 장기 보존 (2년--10년)\\시간 단위 접근, 규제 준수 목적, 최소 비용};

\draw[->, thick, dashed] (hot.south) -- node[right, font=\scriptsize] {90일 경과 시 자동 전환} (warm.north);
\draw[->, thick, dashed] (warm.south) -- node[right, font=\scriptsize] {2년 경과 시 자동 전환} (cold.north);

\end{tikzpicture}
\caption{AI 모델 아카이브를 위한 계층형 스토리지 아키텍처}
\label{fig:tiered_storage}
\end{figure}

\paragraph{비용 최적화 효과}
계층형 스토리지를 적용하면 모든 아카이브를 Hot Tier에 유지하는 대비 \textbf{60--80\%의 스토리지 비용 절감}이 가능하다. 예를 들어 AWS S3 기준, Standard($\sim$\$0.023/GB/월) 대비 Glacier Deep Archive($\sim$\$0.00099/GB/월)는 약 96\%의 비용 차이가 있다. 다만 Cold Tier에서의 복원(Retrieval)에는 추가 비용과 12--48시간의 지연이 발생하므로, 감사 대응이 필요한 아티팩트는 최소 Warm Tier에 유지해야 한다.

\subsection{컴플라이언스 기반 Retention 정책}
\label{sec:12.3.1c}

EU AI Act 제18조(Documentation Keeping)는 고위험 AI 시스템의 기술 문서, 품질 관리 체계 기록, 적합성 선언서를 \textbf{시장 출시 또는 서비스 개시 후 10년간 보존}할 것을 의무화한다\cite{eu2024aiact}. 이는 모델 Retire 이후에도 상당 기간 문서와 아티팩트를 유지해야 함을 의미한다.

\begin{table}[htbp]
\centering
\caption{규제별 Retention 요구사항 비교}
\label{tab:retention_by_regulation}
\begin{tabularx}{\textwidth}{p{3cm}p{2.5cm}p{2.5cm}X}
\toprule
\textbf{규제/표준} & \textbf{대상} & \textbf{Retention Period} & \textbf{보존 범위} \\
\midrule
EU AI Act (제18조) &
고위험 AI 시스템 &
10년 &
기술 문서, 품질 관리 기록, 적합성 선언서, 변경 승인 기록 \\
\midrule
EU AI Act (GPAI) &
범용 AI 모델 &
10년 (최초 출시 후) &
학습 데이터 관련 문서, 저작권 정책, 기술 문서 \\
\midrule
GDPR (제30조) &
개인정보 처리 &
처리 목적 달성 시까지 &
처리 활동 기록, 동의 기록, 영향 평가 \\
\midrule
금융 규제 (바젤) &
금융 AI 모델 &
5--7년 &
모델 검증 보고서, 성능 기록, 의사결정 로그 \\
\midrule
의료기기 규제 &
AI 의료기기 &
기기 수명 + 15년 &
임상 평가, 설계 이력, 소프트웨어 변경 기록 \\
\midrule
ISO/IEC 42001 &
AI 관리 체계 &
인증 주기 (3년) + $\alpha$ &
AI 정책, 리스크 평가, 성과 평가 기록 \\
\bottomrule
\end{tabularx}
\end{table}

\paragraph{Retention 정책 설계 원칙}
조직의 Retention 정책을 설계할 때 다음 원칙을 적용한다.

\begin{enumerate}
 \item \textbf{최장 의무 기준 적용}: 복수 규제가 적용되는 경우 가장 긴 Retention Period를 기준으로 한다. 예를 들어 EU AI Act 고위험 시스템이면서 의료기기인 경우, 의료기기 규제의 15년+ 기준을 적용한다.
 \item \textbf{자동 만료 및 삭제}: Retention Period 경과 후 자동 삭제 파이프라인을 구축하여 불필요한 데이터 보유에 따른 리스크를 최소화한다. 특히 GDPR의 데이터 최소화 원칙과의 균형이 중요하다.
 \item \textbf{법적 보존(Legal Hold) 메커니즘}: 소송, 규제 조사 등 예외 상황에서 자동 삭제를 중단하는 Legal Hold 기능을 구현한다.
 \item \textbf{보존 증명(Chain of Custody)}: 아카이브된 아티팩트의 무결성을 증명할 수 있도록 해시값, 타임스탬프, 접근 로그를 유지한다.
\end{enumerate}

\subsection{기계 망각(Machine Unlearning)과 보존 정책}
\label{sec:12.3.2}

Privacy 규정이 강화됨에 따라 Archiving 정책에는 'Data 삭제 요청'에 대한 Response 방안이 포함되어야 한다. 특정 데이터 주체의 삭제 요청 시, 모델 아카이브 내 가중치에서 해당 정보를 제거하는 \textbf{머신 언러닝(Machine Unlearning)} \cite{cao2015towards} Technical을 적용하거나, 해당 버전의 아카이브를 Retire하는 정책을 수립한다.

\paragraph{SISA 프레임워크 개요}
SISA(Sharded, Isolated, Sliced, Aggregated)\cite{bourtoule2021unlearning}는 머신 언러닝을 위한 대표적 프레임워크로, 학습 데이터를 사전에 분할(Shard)하고 각 샤드 내에서 점진적으로 슬라이스(Slice)를 추가하며 학습한다. 특정 데이터 포인트의 삭제 요청 시 해당 포인트가 속한 샤드의 구성 모델만 재학습하면 되므로, 전체 모델 재학습 대비 \textbf{계산 비용을 95\% 이상 절감}할 수 있다.

\begin{figure}[htbp]
\centering
\begin{tikzpicture}[
 shard/.style={rectangle, draw, minimum width=2cm, minimum height=1.5cm, align=center, font=\small},
 slice/.style={rectangle, draw, minimum width=1.8cm, minimum height=0.3cm, align=center, font=\scriptsize},
 model/.style={circle, draw, fill=green!20, minimum size=1cm, align=center, font=\scriptsize},
 agg/.style={rectangle, draw, fill=red!20, rounded corners, minimum width=2.5cm, minimum height=1cm, align=center, font=\small},
]

% Shards
\node[shard, fill=blue!10] (s1) at (0, 0) {};
\node[above, font=\small\bfseries] at (s1.north) {Shard 1};
\node[slice, fill=blue!20] at (0, 0.3) {Slice 1-1};
\node[slice, fill=blue!30] at (0, 0) {Slice 1-2};
\node[slice, fill=blue!40] at (0, -0.3) {Slice 1-3};

\node[shard, fill=orange!10] (s2) at (3, 0) {};
\node[above, font=\small\bfseries] at (s2.north) {Shard 2};
\node[slice, fill=orange!20] at (3, 0.3) {Slice 2-1};
\node[slice, fill=orange!30] at (3, 0) {Slice 2-2};
\node[slice, fill=orange!40] at (3, -0.3) {Slice 2-3};

\node[shard, fill=green!10] (s3) at (6, 0) {};
\node[above, font=\small\bfseries] at (s3.north) {Shard 3};
\node[slice, fill=green!20] at (6, 0.3) {Slice 3-1};
\node[slice, fill=green!30] at (6, 0) {Slice 3-2};
\node[slice, fill=green!40] at (6, -0.3) {Slice 3-3};

% Constituent Models
\node[model] (m1) at (0, -2.2) {$M_1$};
\node[model] (m2) at (3, -2.2) {$M_2$};
\node[model] (m3) at (6, -2.2) {$M_3$};

\draw[->, thick] (s1) -- (m1);
\draw[->, thick] (s2) -- (m2);
\draw[->, thick] (s3) -- (m3);

% Aggregation
\node[agg] (agg) at (3, -4) {\textbf{Aggregation}\\(다수결/평균)};
\draw[->, thick] (m1) -- (agg);
\draw[->, thick] (m2) -- (agg);
\draw[->, thick] (m3) -- (agg);

\end{tikzpicture}
\caption{SISA 프레임워크 구조: 데이터를 Shard/Slice로 분할하여 학습}
\label{fig:sisa_architecture}
\end{figure}

다음 코드는 SISA 기반 머신 언러닝 시스템의 핵심 구현을 보여준다.

\begin{lstlisting}[language=Python, caption={SISA 기반 머신 언러닝 구현}, label={lst:sisa_unlearning}]
"""
SISA (Sharded, Isolated, Sliced, Aggregated) Machine Unlearning
- 학습 데이터를 Shard로 분할, 각 Shard 내에서 Slice 단위 점진 학습
- 삭제 요청 시 해당 Shard의 구성 모델만 재학습
- Bourtoule et al., 2021 기반 구현
"""

from dataclasses import dataclass, field
from typing import List, Dict, Optional, Any, Callable
from datetime import datetime
import hashlib
import numpy as np
import logging

logger = logging.getLogger(__name__)

@dataclass
class UnlearningRequest:
    """데이터 삭제(언러닝) 요청"""
    request_id: str
    data_subject_id: str
    data_point_ids: List[str]
    reason: str  # "GDPR Art.17", "CCPA", "voluntary"
    requested_at: datetime = field(
        default_factory=datetime.now)
    completed_at: Optional[datetime] = None
    status: str = "pending"  # pending, processing, completed

@dataclass
class ShardCheckpoint:
    """Shard별 학습 체크포인트"""
    shard_id: int
    slice_index: int
    model_weights: Any  # 직렬화된 모델 가중치
    training_data_ids: List[str]
    created_at: datetime = field(
        default_factory=datetime.now)
    checksum: str = ""

    def compute_checksum(self) -> str:
        """체크포인트 무결성 해시"""
        content = f"{self.shard_id}:{self.slice_index}"
        content += f":{len(self.training_data_ids)}"
        self.checksum = hashlib.sha256(
            content.encode()).hexdigest()[:16]
        return self.checksum

class SISATrainer:
    """
    SISA 학습 및 언러닝 관리자

    핵심 아이디어:
    - 전체 학습 데이터를 num_shards개 샤드로 분할
    - 각 샤드를 num_slices개 슬라이스로 나누어 점진 학습
    - 각 슬라이스 학습 완료 시 체크포인트 저장
    - 삭제 요청 시, 해당 데이터가 속한 샤드만 재학습
    """

    def __init__(
        self,
        num_shards: int = 5,
        num_slices: int = 10,
        model_factory: Optional[Callable] = None,
        aggregation: str = "majority_vote",
    ):
        self.num_shards = num_shards
        self.num_slices = num_slices
        self.model_factory = model_factory
        self.aggregation = aggregation

        # 샤드별 구성 모델
        self.constituent_models: Dict[int, Any] = {}
        # 데이터 포인트 -> 샤드 매핑
        self.data_to_shard: Dict[str, int] = {}
        # 샤드별 체크포인트 히스토리
        self.checkpoints: Dict[int, List[ShardCheckpoint]] = {
            i: [] for i in range(num_shards)
        }
        # 삭제 요청 이력
        self.unlearn_history: List[UnlearningRequest] = []

    def assign_data_to_shards(
        self,
        data_ids: List[str],
    ) -> Dict[int, List[str]]:
        """데이터를 샤드에 균등 배정 (해시 기반)"""
        shard_assignments: Dict[int, List[str]] = {
            i: [] for i in range(self.num_shards)
        }

        for data_id in data_ids:
            # 결정론적 샤드 배정 (재현 가능)
            hash_val = int(hashlib.md5(
                data_id.encode()).hexdigest(), 16)
            shard_id = hash_val % self.num_shards
            shard_assignments[shard_id].append(data_id)
            self.data_to_shard[data_id] = shard_id

        logger.info(
            f"데이터 {len(data_ids)}건을 "
            f"{self.num_shards}개 샤드에 배정 완료"
        )
        return shard_assignments

    def train_shard(
        self,
        shard_id: int,
        shard_data: List[str],
        from_slice: int = 0,
    ) -> None:
        """
        특정 샤드의 구성 모델 학습
        from_slice > 0이면 해당 슬라이스의 체크포인트부터 재개
        """
        slice_size = max(1, len(shard_data) // self.num_slices)

        # 체크포인트에서 재개
        if from_slice > 0 and self.checkpoints[shard_id]:
            checkpoint = self.checkpoints[shard_id][
                from_slice - 1]
            model = self._load_model(checkpoint.model_weights)
            logger.info(
                f"Shard {shard_id}: "
                f"Slice {from_slice}부터 재개")
        else:
            model = self.model_factory()
            from_slice = 0

        # 슬라이스 단위 점진 학습
        for s in range(from_slice, self.num_slices):
            start_idx = 0
            end_idx = min((s + 1) * slice_size,
                          len(shard_data))
            current_data = shard_data[:end_idx]

            # 모델 학습 (구현체에 위임)
            model = self._train_model(
                model, current_data)

            # 체크포인트 저장
            checkpoint = ShardCheckpoint(
                shard_id=shard_id,
                slice_index=s,
                model_weights=self._serialize_model(model),
                training_data_ids=current_data.copy(),
            )
            checkpoint.compute_checksum()

            # 기존 체크포인트 교체 또는 추가
            if s < len(self.checkpoints[shard_id]):
                self.checkpoints[shard_id][s] = checkpoint
            else:
                self.checkpoints[shard_id].append(
                    checkpoint)

            logger.info(
                f"Shard {shard_id}, Slice {s} "
                f"학습 완료 (데이터: {len(current_data)}건)")

        self.constituent_models[shard_id] = model

    def process_unlearn_request(
        self,
        request: UnlearningRequest,
        full_shard_data: Dict[int, List[str]],
    ) -> Dict:
        """
        언러닝 요청 처리
        1. 삭제 대상 데이터가 속한 샤드 식별
        2. 해당 데이터 제거 후 영향받는 샤드만 재학습
        3. 재학습은 가장 가까운 체크포인트에서 시작
        """
        request.status = "processing"
        affected_shards = set()
        retrain_info = {}

        # 1. 영향받는 샤드 식별
        for data_id in request.data_point_ids:
            shard_id = self.data_to_shard.get(data_id)
            if shard_id is not None:
                affected_shards.add(shard_id)

        logger.info(
            f"언러닝 요청 {request.request_id}: "
            f"{len(request.data_point_ids)}건 삭제, "
            f"{len(affected_shards)}개 샤드 영향")

        # 2. 영향받는 샤드만 재학습
        for shard_id in affected_shards:
            # 삭제 대상 제거
            original_data = full_shard_data[shard_id]
            cleaned_data = [
                d for d in original_data
                if d not in request.data_point_ids
            ]

            # 가장 가까운 안전 체크포인트 탐색
            safe_slice = self._find_safe_checkpoint(
                shard_id, request.data_point_ids)

            # 재학습 실행
            self.train_shard(
                shard_id, cleaned_data,
                from_slice=safe_slice)

            retrain_info[shard_id] = {
                'original_size': len(original_data),
                'cleaned_size': len(cleaned_data),
                'retrain_from_slice': safe_slice,
            }

            # 매핑 정보 업데이트
            for data_id in request.data_point_ids:
                self.data_to_shard.pop(data_id, None)

            full_shard_data[shard_id] = cleaned_data

        # 3. 완료 기록
        request.status = "completed"
        request.completed_at = datetime.now()
        self.unlearn_history.append(request)

        return {
            'request_id': request.request_id,
            'affected_shards': list(affected_shards),
            'retrain_info': retrain_info,
            'total_shards': self.num_shards,
            'unaffected_shards': (
                self.num_shards - len(affected_shards)),
        }

    def _find_safe_checkpoint(
        self,
        shard_id: int,
        removed_ids: List[str],
    ) -> int:
        """삭제 대상이 포함되지 않은 가장 마지막 체크포인트"""
        removed_set = set(removed_ids)

        for i, cp in enumerate(
                self.checkpoints[shard_id]):
            if removed_set & set(cp.training_data_ids):
                return max(0, i)

        return 0

    def _train_model(self, model, data):
        """모델 학습 (구현체에 위임)"""
        # 실제 구현에서는 PyTorch/TensorFlow 학습 루프
        return model

    def _serialize_model(self, model) -> bytes:
        """모델 직렬화"""
        return b""  # 실제: pickle/torch.save

    def _load_model(self, weights: bytes):
        """체크포인트에서 모델 복원"""
        return self.model_factory()

    def get_unlearning_audit_log(self) -> List[Dict]:
        """언러닝 감사 로그 생성 (규제 대응)"""
        return [
            {
                'request_id': r.request_id,
                'data_subject': r.data_subject_id,
                'num_points_removed': len(r.data_point_ids),
                'reason': r.reason,
                'requested_at': r.requested_at.isoformat(),
                'completed_at': (
                    r.completed_at.isoformat()
                    if r.completed_at else None),
                'status': r.status,
            }
            for r in self.unlearn_history
        ]
\end{lstlisting}

\paragraph{머신 언러닝과 아카이브 정책의 조화}
머신 언러닝 요청이 아카이브된 모델에 대해 발생하는 경우, 다음 세 가지 대응 전략 중 하나를 선택한다.

\begin{enumerate}
 \item \textbf{아카이브 모델 재학습}: SISA 체크포인트를 활용하여 해당 데이터를 제거한 버전으로 재학습 후 아카이브를 교체한다. 비용이 발생하지만 가장 완전한 언러닝을 보장한다.
 \item \textbf{아카이브 삭제}: 해당 버전의 모델 아카이브 자체를 삭제하고, 삭제 사실과 사유를 메타데이터로 기록한다. 규제 보존 의무와 충돌할 수 있으므로 법무팀 협의가 필요하다.
 \item \textbf{접근 제한 및 라벨링}: 아카이브를 유지하되 ``언러닝 미완료'' 라벨을 부착하고, 해당 모델의 추론 결과 사용을 금지한다. 감사 목적의 보존만 허용한다.
\end{enumerate}


%----------------------------------------------------------
\section{모델 Retire와 생성형 AI}
\label{sec:12.3b}

2023년 이후 대규모 언어 모델(LLM) 기반 시스템의 엔터프라이즈 도입이 급증하면서, 기존 전통적 ML 모델과는 본질적으로 다른 Retire 과제가 등장하였다. LLM 파인튜닝 모델, RAG(Retrieval-Augmented Generation) 시스템, 프롬프트\cite{openai2023gpt4} 템플릿 등 생성형 AI 고유의 아티팩트에 대한 체계적인 Retire 전략이 필요하다.

\subsection{LLM 파인튜닝 모델 Retire}
\label{sec:12.3b.1}

LLM 파인튜닝 모델의 Retire는 다음과 같은 고유한 복잡성을 가진다.

\begin{itemize}
 \item \textbf{기반 모델 의존성}: 파인튜닝 모델은 기반 모델(Base Model)에 종속된다. 기반 모델 제공사(OpenAI, Anthropic, Google 등)가 해당 버전을 Deprecate하면 파인튜닝 모델도 강제 Retire된다. Azure AI Foundry, OpenAI API 등은 공식 Deprecation 일정을 사전 공지하며, 통상 6개월의 Migration 기간을 제공한다.
 \item \textbf{어댑터 가중치 관리}: LoRA, QLoRA 등 파라미터 효율적 파인튜닝(PEFT) 방식으로 학습된 어댑터 가중치는 기반 모델 버전에 긴밀하게 결합된다. 기반 모델 버전이 변경되면 어댑터의 호환성이 보장되지 않으므로, 아카이브 시 기반 모델 버전 정보를 반드시 함께 기록해야 한다.
 \item \textbf{학습 데이터 규모와 비용}: LLM 파인튜닝에 사용된 데이터셋은 수만에서 수십만 건에 달할 수 있으며, 이 데이터의 큐레이션 비용이 상당하다. Retire 시 학습 데이터와 데이터 파이프라인을 함께 아카이브하여 향후 재활용 가능성을 보존한다.
 \item \textbf{평가 벤치마크}: 파인튜닝 모델의 성능을 측정한 도메인별 평가 데이터셋과 벤치마크 결과를 아카이브하여, 후속 모델의 성능 비교 기준선으로 활용한다.
\end{itemize}

\begin{table}[htbp]
\centering
\caption{LLM 파인튜닝 모델 Retire 시 아카이브 항목}
\label{tab:llm_finetune_archive}
\begin{tabularx}{\textwidth}{p{3cm}p{3.5cm}X}
\toprule
\textbf{아티팩트} & \textbf{상세 내용} & \textbf{보존 근거} \\
\midrule
어댑터 가중치 &
LoRA/QLoRA 어댑터 파일, 설정 &
재현성, 롤백 대비 \\
\midrule
기반 모델 정보 &
모델명, 버전, 파라미터 수, 체크포인트 해시 &
호환성 추적, 의존성 관리 \\
\midrule
학습 데이터셋 &
원본 데이터 참조, 전처리 파이프라인, 데이터 카드 &
재학습 가능성, 지식 보존 \\
\midrule
학습 설정 &
하이퍼파라미터, 학습률 스케줄, 에폭 수 &
재현성 \\
\midrule
평가 결과 &
벤치마크 점수, 사람 평가 결과, A/B 테스트 결과 &
성능 비교 기준선 \\
\midrule
안전성 평가 &
Red teaming 결과, 편향 테스트, 유해성 평가 &
규제 준수, 감사 대응 \\
\bottomrule
\end{tabularx}
\end{table}

\subsection{RAG 시스템 구성요소 Retire}
\label{sec:12.3b.2}

RAG 시스템은 LLM, 벡터 데이터베이스, 임베딩 모델, 문서 파이프라인 등 다수의 구성요소가 결합된 복합 시스템이다. 한 구성요소의 Retire가 전체 시스템에 연쇄적 영향을 미칠 수 있으므로, 구성요소 간 의존성을 파악한 후 Retire 순서를 결정해야 한다.

\begin{figure}[htbp]
\centering
\begin{tikzpicture}[
 comp/.style={rectangle, draw, rounded corners, minimum width=2.5cm, minimum height=0.9cm, align=center, font=\small},
 arrow/.style={->, thick},
]

% 구성요소
\node[comp, fill=blue!20] (doc) at (0, 2) {문서 소스\\(Knowledge Base)};
\node[comp, fill=green!20] (pipe) at (4, 2) {청킹/전처리\\파이프라인};
\node[comp, fill=orange!20] (embed) at (8, 2) {임베딩 모델};
\node[comp, fill=purple!20] (vector) at (4, 0) {벡터\\데이터베이스};
\node[comp, fill=red!20] (llm) at (8, 0) {LLM\\(생성 모델)};
\node[comp, fill=yellow!20] (prompt) at (0, 0) {프롬프트\\템플릿};

% 화살표
\draw[arrow] (doc) -- (pipe);
\draw[arrow] (pipe) -- (embed);
\draw[arrow] (embed) -- (vector);
\draw[arrow] (vector) -- (llm);
\draw[arrow] (prompt) -- (llm);

\end{tikzpicture}
\caption{RAG 시스템 구성요소 의존성 맵}
\label{fig:rag_dependency}
\end{figure}

\paragraph{구성요소별 Retire 고려사항}

\begin{itemize}
 \item \textbf{벡터 데이터베이스}: 임베딩 모델이 변경되면 기존 벡터 인덱스는 무용화된다. 임베딩 모델 Retire 시 벡터 DB의 전체 재인덱싱(Re-indexing) 또는 동시 Retire를 계획한다.
 \item \textbf{임베딩 모델}: 동일 임베딩 모델을 여러 RAG 시스템이 공유하는 경우, Retire 영향 범위가 확대된다. 제9장의 리니지 시스템을 활용하여 다운스트림 소비자를 사전 파악한다.
 \item \textbf{문서 소스}: Knowledge Base의 폐기는 해당 문서에 대한 RAG 응답 품질에 직접적 영향을 미친다. 문서 소스 Retire 전 대체 소스 확보 또는 시스템 동시 Retire를 결정한다.
 \item \textbf{프롬프트 템플릿}: 프롬프트의 변경/폐기가 LLM 출력 품질에 미치는 영향을 사전 평가한다. 특히 Few-shot 예제가 포함된 프롬프트는 도메인 지식의 손실을 의미할 수 있다.
\end{itemize}

\subsection{프롬프트 템플릿 버저닝과 Archiving}
\label{sec:12.3b.3}

생성형 AI 시대에서 프롬프트는 모델의 행동을 결정하는 핵심 아티팩트이며, 소프트웨어 코드와 동등한 수준의 버전 관리가 필요하다.

\begin{table}[htbp]
\centering
\caption{프롬프트 템플릿 Archiving 기준}
\label{tab:prompt_archiving}
\begin{tabularx}{\textwidth}{p{3cm}p{2.5cm}X}
\toprule
\textbf{아티팩트} & \textbf{Retention Period} & \textbf{보존 사유 및 방법} \\
\midrule
시스템 프롬프트 &
모델 Retire + 5년 &
모델 행동 재현, 감사 대응; Git 기반 버전 관리 \\
\midrule
Few-shot 예제 &
모델 Retire + 3년 &
도메인 지식 보존; 데이터 카탈로그 연계 \\
\midrule
프롬프트 평가 결과 &
모델 Retire + 3년 &
품질 비교 기준; 벤치마크 DB 저장 \\
\midrule
가드레일 규칙 &
Permanent &
안전성 정책 이력; 정책 관리 시스템 연계 \\
\midrule
A/B 테스트 결과 &
2년 &
프롬프트 최적화 지식 보존 \\
\bottomrule
\end{tabularx}
\end{table}

\paragraph{프롬프트 거버넌스 원칙}
프롬프트 템플릿의 Retire와 Archiving에 다음 원칙을 적용한다.

\begin{enumerate}
 \item \textbf{코드로서의 프롬프트(Prompt-as-Code)}: 모든 프롬프트를 Git 리포지토리에서 관리하며, 변경 이력, 리뷰 기록, 승인 내역을 추적한다.
 \item \textbf{프롬프트-모델 결합 기록}: 특정 프롬프트가 어떤 모델 버전과 함께 사용되었는지를 메타데이터로 기록한다. 모델이 Retire되더라도 프롬프트 아카이브를 통해 당시의 시스템 행동을 이해할 수 있어야 한다.
 \item \textbf{민감 정보 스크러빙}: 프롬프트 내 개인정보, 비밀 정보(API 키 등)가 포함되지 않았는지 아카이브 전 자동 점검한다.
 \item \textbf{크로스 모델 재사용성 평가}: Retire되는 모델의 프롬프트 자산 중 후속 모델에서 재사용 가능한 항목을 식별하여 프롬프트 라이브러리로 이관한다.
\end{enumerate}


%----------------------------------------------------------
\section{지식 보존과 이전}
\label{sec:12.4}

\subsection{Retire Retrospective(Retrospective)}
\label{sec:12.4.2}

Model Retire는 단순히 시스템을 끄는 행위가 아니라 하나의 실험적 결말이다. 왜 실패했는지, 어떤 시도가 무의미했는지를 기록하여 조직의 자산으로 만든다.

\begin{lstlisting}[language=Python, caption={Retirement Retrospective and Lessons Learned Management System}, label={lst:retirement_retrospective}]
class LessonCategory(Enum):
 DATA = "Data"; MODELING = "Modeling"; GOVERNANCE = "Governance"

@dataclass
class LessonLearned:
 lesson_id: str; category: LessonCategory; title: str
 description: str; context: str; recommendation: str

class RetirementRetrospectiveManager:
 """Manager of turning mistakes into knowledge"""
 def add_lesson(self, retro_id: str, lesson: LessonLearned):
 # Register lessons in Knowledge Base (KB) for future retrieval
 pass
\end{lstlisting}

\subsection{Technical 부채 관리와 Retire의 연계}
\label{sec:12.4.3}

ML 시스템의 Technical 부채(Technical Debt)는 Sculley et al.\cite{sculley2015hidden}이 지적한 것처럼 전통적 소프트웨어보다 더욱 빠르게 누적된다. 모델 Retire는 이러한 Technical 부채를 체계적으로 청산하는 핵심 메커니즘이다.

\paragraph{AI Technical 부채의 유형}

\begin{table}[htbp]
\centering
\caption{AI Technical 부채 유형과 Retire를 통한 해소}
\label{tab:tech_debt_types}
\begin{tabularx}{\textwidth}{p{2.5cm}p{4cm}X}
\toprule
\textbf{부채 유형} & \textbf{설명} & \textbf{Retire 시 해소 방안} \\
\midrule
데이터 의존성 부채 &
불안정한 데이터 소스에 대한 암묵적 의존; 미사용 피처의 누적 &
데이터 리니지 정리; 미사용 파이프라인 제거; 피처 스토어 정리 \\
\midrule
설정 부채 &
수작업으로 관리되는 하이퍼파라미터; 환경별 설정 파편화 &
설정 아카이브 후 표준화된 구성 관리로 전환 \\
\midrule
모니터링 부채 &
불충분한 모니터링으로 인한 ``Silent Failure'' 리스크 &
모니터링 대상 축소; 알림 규칙 정리; 대시보드 간소화 \\
\midrule
피드백 루프 부채 &
모델 출력이 미래 학습 데이터에 영향을 미치는 숨겨진 피드백 루프 &
피드백 루프 차단; 영향받은 다운스트림 모델 재평가 \\
\midrule
코드 부채 &
실험용 코드의 프로덕션 잔류; 글루 코드(Glue Code) 과다 &
레거시 코드 제거; Integration 포인트 정리 \\
\bottomrule
\end{tabularx}
\end{table}

\paragraph{Retire를 통한 부채 청산 프로세스}
모델 Retire 시 다음 절차를 통해 Technical 부채를 체계적으로 청산한다.

\begin{enumerate}
 \item \textbf{의존성 맵 작성}: Retire 대상 모델이 참조하는 데이터 소스, 피처, 외부 서비스를 모두 식별한다.
 \item \textbf{공유 리소스 분석}: 해당 모델만 사용하는 리소스(전용)와 다른 모델과 공유하는 리소스를 구분한다.
 \item \textbf{전용 리소스 회수}: Retire 모델 전용의 데이터 파이프라인, 피처 변환 코드, 모니터링 설정을 제거한다.
 \item \textbf{공유 리소스 재평가}: 공유 리소스의 경우, Retire 후 남은 소비자가 실제로 해당 리소스를 사용하는지 검증한다.
 \item \textbf{부채 청산 기록}: 제거된 코드, 파이프라인, 설정의 목록과 절감된 복잡도를 기록한다.
\end{enumerate}


%----------------------------------------------------------
\section{사후 영향 분석 (제9장 데이터 계보 연계)}
\label{sec:12.5}

\subsection{계보-based Impact Analysis}
\label{sec:12.5.1}

제9장에서 구축한 통합 리니지 시스템(통합된 계보 시스템, 9.3절 참조)을 활용하여 모델 폐기가 다운스트림 시스템에 미치는 영향을 분석한다. 특히 데이터 패브릭의 메타데이터를 활용하여, 원천 데이터의 보존 기한 만료 시 이를 학습한 모델에 '자동 폐기 알림'을 보내는 체계를 권고한다. 이는 9장에서 설명한 통합 가시성을 실제 운영 효율화(ROI)로 연결하는 핵심 사례이다.

\begin{lstlisting}[language=Python, caption={Lineage-based Retirement Impact Analyzer}, label={lst:retirement_impact_analysis}]
class DependencyType(Enum):
 DIRECT_CONSUMER = "Direct Consumer"; ENSEMBLE_MEMBER = "Ensemble Member"
 FALLBACK_MODEL = "Fallback Model"

class RetirementImpactAnalyzer:
 """Cascade failure analysis via lineage graph"""
 def analyze_impact(self, model_id: str, lineage_manager):
 # Call IntegratedLineageManager from Ch 9 for downstream tracking
 downstream = lineage_manager.get_downstream(model_id)
 # Impact Retire (Blocker) Setup
 pass
\end{lstlisting}


%----------------------------------------------------------
\section{실무 도구}
\label{sec:12.6}

\begin{practicaltool}{AI 운영 지표 대시보드 설계 가이드}
\textbf{[실무 도구 12-1]} 다음 가이드를 활용하여 AI 모델 운영 및 Retire 판단을 위한 대시보드를 설계한다.

\paragraph{대시보드 구성 원칙}
AI 운영 지표 대시보드는 다음 4개 영역으로 구성하며, 각 영역은 모델 Retire 의사결정에 필요한 핵심 지표를 실시간으로 제공한다.

\begin{enumerate}
 \item \textbf{모델 건강도 패널 (Model Health Panel)}
 \begin{itemize}
  \item 정확도/F1/AUC 추이 그래프 (기준선 대비 현재값)
  \item 데이터 드리프트 점수 (PSI, KL Divergence)
  \item 예측 분포 변화 히스토그램
  \item 자동 알림: 성능 지표가 임계값 하회 시 Slack/Teams 알림
 \end{itemize}

 \item \textbf{비즈니스 가치 패널 (Business Value Panel)}
 \begin{itemize}
  \item 일/주/월 별 예측 요청 수 추이
  \item 활성 소비자(Consumer) 수 추이
  \item 월간 추정 가치 vs 월간 운영 비용 (ROI 시각화)
  \item 비용 효율성 순위 (전사 모델 대비)
 \end{itemize}

 \item \textbf{운영 안정성 패널 (Operational Stability Panel)}
 \begin{itemize}
  \item 평균 응답 시간(Latency) 추이
  \item 에러율 및 타임아웃 비율
  \item 최근 90일 인시던트 카운트
  \item 가용성(Uptime) 비율
 \end{itemize}

 \item \textbf{거버넌스 컴플라이언스 패널 (Governance Panel)}
 \begin{itemize}
  \item 최종 감사일 및 다음 예정 감사일
  \item 미해결 감사 지적사항 수
  \item 문서 완성도 비율 (Model Card, 평가서 등)
  \item 프레임워크/라이브러리 EOL 경고
  \item 보안 취약점(CVE) 카운트
 \end{itemize}
\end{enumerate}

\paragraph{좀비 모델 경고 영역}
대시보드 상단에 \textbf{좀비 모델 경고 배너}를 배치한다.

\begin{itemize}
 \item Critical/High 등급 좀비 모델 수와 예상 월간 낭비 비용을 표시
 \item 클릭 시 좀비 탐지 상세 리포트로 이동
 \item 주간 스캔 결과를 자동 갱신
\end{itemize}

\paragraph{기술 스택 참조}
\begin{itemize}
 \item \textbf{데이터 수집}: Prometheus + Grafana 또는 Datadog
 \item \textbf{드리프트 모니터링}: Evidently AI, WhyLabs, NannyML
 \item \textbf{대시보드}: Grafana 대시보드 또는 Streamlit 커스텀 대시보드
 \item \textbf{알림}: PagerDuty, Slack Webhook, Microsoft Teams Connector
 \item \textbf{모델 레지스트리 연동}: MLflow, AWS SageMaker Model Registry, Azure ML
\end{itemize}

\paragraph{갱신 주기}
\begin{itemize}
 \item 건강도/운영 지표: 실시간 또는 5분 간격
 \item 비즈니스 가치: 일간 집계
 \item 거버넌스 지표: 주간 또는 감사 이벤트 발생 시
 \item 좀비 탐지: 주간 스캔 결과 반영
\end{itemize}
\end{practicaltool}


\begin{practicaltool}{모델 폐기(Retire) 체크리스트}
\textbf{[실무 도구 12-2]} 다음 체크리스트를 활용하여 모델 Retire를 체계적으로 수행한다.

\paragraph{A. Retire 사전 준비}
\begin{enumerate}
 \item[$\square$] Retire 트리거 확인 및 문서화 완료
 \item[$\square$] 모델 건강도 평가(12.1.3절 프레임워크) 수행 완료
 \item[$\square$] 영향받는 다운스트림 시스템 목록 작성 완료 (리니지 기반)
 \item[$\square$] 영향받는 이해관계자 목록 작성 완료
 \item[$\square$] 대체 모델/시스템 준비 상태 확인
 \item[$\square$] Retire 일정 및 마일스톤 수립
\end{enumerate}

\paragraph{B. 승인 및 통보}
\begin{enumerate}
 \item[$\square$] AI 거버넌스 위원회 Retire 승인 획득
 \item[$\square$] CDO/CTO 최종 승인
 \item[$\square$] 이해관계자 사전 통보 (Planned: 30일 전 / Urgent: 7일 전)
 \item[$\square$] 소비자 시스템 담당자와 Migration 일정 합의
 \item[$\square$] 고객 대면 서비스인 경우 고객 공지
\end{enumerate}

\paragraph{C. 아카이브 수행}
\begin{enumerate}
 \item[$\square$] 모델 가중치(최종 버전 + 주요 이전 버전) 아카이브
 \item[$\square$] 모델 코드(학습/추론/전처리) 아카이브
 \item[$\square$] 학습 데이터 참조(위치, 버전, 스키마) 아카이브
 \item[$\square$] 하이퍼파라미터 및 학습 설정 아카이브
 \item[$\square$] Model Card 및 평가 보고서 아카이브
 \item[$\square$] 승인 기록 및 감사 이력 아카이브
 \item[$\square$] (LLM) 어댑터 가중치 + 기반 모델 정보 아카이브
 \item[$\square$] (RAG) 벡터 DB 스냅샷 또는 메타데이터 아카이브
 \item[$\square$] (생성형 AI) 프롬프트 템플릿 및 가드레일 규칙 아카이브
 \item[$\square$] 아카이브 무결성 검증(체크섬 확인)
 \item[$\square$] Retention Period 설정 및 자동 계층 전환 정책 적용
\end{enumerate}

\paragraph{D. 서비스 종료 및 리소스 회수}
\begin{enumerate}
 \item[$\square$] 트래픽 대체 시스템으로 100\% 전환 확인
 \item[$\square$] API 엔드포인트 비활성화
 \item[$\square$] 모니터링 알림 비활성화
 \item[$\square$] 컴퓨팅 리소스(GPU/CPU 인스턴스) 회수
 \item[$\square$] 스토리지 리소스 정리(Hot Tier에서 제거)
 \item[$\square$] IAM 권한 및 서비스 계정 비활성화
 \item[$\square$] CI/CD 파이프라인 비활성화
\end{enumerate}

\paragraph{E. 마무리 및 기록}
\begin{enumerate}
 \item[$\square$] 모델 레지스트리 상태를 ``Retired''로 갱신
 \item[$\square$] Retire Retrospective 수행 및 Lessons Learned 기록
 \item[$\square$] Technical 부채 청산 내역 기록 (제거된 코드, 파이프라인, 설정)
 \item[$\square$] 비용 절감 효과 산출 및 보고
 \item[$\square$] Retire 완료 보고서 작성 및 거버넌스 위원회 보고
 \item[$\square$] 관련 문서 최종 갱신 (아키텍처 다이어그램, 시스템 목록 등)
\end{enumerate}

\paragraph{F. 컴플라이언스 확인}
\begin{enumerate}
 \item[$\square$] EU AI Act 10년 보존 의무 대상 여부 확인
 \item[$\square$] GDPR 데이터 최소화 원칙과 보존 의무의 균형 검토
 \item[$\square$] 산업별 추가 규제 요건 확인 (금융, 의료 등)
 \item[$\square$] 미결 머신 언러닝 요청 처리 완료 확인
 \item[$\square$] Legal Hold 대상 여부 확인
\end{enumerate}
\end{practicaltool}


%----------------------------------------------------------
\subsection*{제12장 요약}

본 장에서는 AI 모델 Retire와 Archiving의 전 과정을 체계적으로 다루었다. 핵심 내용을 정리하면 다음과 같다.

\begin{table}[htbp]
\centering
\caption{제12장 핵심 요약}
\label{tab:ch12_summary}
\begin{tabularx}{\textwidth}{p{3.5cm}X}
\toprule
\textbf{주제} & \textbf{핵심 내용} \\
\midrule
Retire 기준과 트리거 &
Technical, Business, 규제, Operational의 4대 요인으로 분류. 정량적 건강도 평가 프레임워크를 통해 객관적 Retire 의사결정 지원. \\
\midrule
좀비 모델 탐지 &
사용량, 소유자, 비용, 프레임워크 EOL을 기반으로 한 자동 탐지 시스템 구현. 주간/월간 정기 스캔으로 거버넌스 사각지대 해소. \\
\midrule
안전한 Retire 프로세스 &
9단계 표준 절차(평가 $\rightarrow$ 승인 $\rightarrow$ 통보 $\rightarrow$ Migration $\rightarrow$ Archiving $\rightarrow$ Shutdown $\rightarrow$ 회수 $\rightarrow$ 문서화 $\rightarrow$ 완료). 긴급도별 차등 절차 운영. \\
\midrule
Archiving 전략 &
계층형 스토리지(Hot/Warm/Cold)를 통한 비용 최적화. 클라우드 vs 온프레미스 하이브리드 전략. EU AI Act 10년 보존 의무 대응. \\
\midrule
머신 언러닝 &
SISA 프레임워크를 통한 효율적 데이터 삭제. GDPR 삭제 요청과 규제 보존 의무 간 균형 전략. 감사 로그 유지. \\
\midrule
생성형 AI Retire &
LLM 파인튜닝 모델의 기반 모델 의존성 관리. RAG 시스템 구성요소 간 연쇄 Retire. 프롬프트 템플릿 버저닝과 Archiving. \\
\midrule
지식 보존 &
Retire Retrospective를 통한 조직 학습. Technical 부채 청산과 Retire의 연계. 리니지 기반 사후 영향 분석. \\
\midrule
실무 도구 &
AI 운영 지표 대시보드 설계 가이드(실무 도구 12-1). 모델 폐기 체크리스트(실무 도구 12-2). \\
\bottomrule
\end{tabularx}
\end{table}

모델 Retire는 AI 거버넌스 생명주기의 마지막 단계이자, 동시에 다음 세대 모델을 위한 출발점이다. 체계적인 Retire 프로세스를 통해 조직은 (1) 좀비 모델에 의한 불필요한 비용과 리스크를 제거하고, (2) 규제 요건에 대한 선제적 준수를 확보하며, (3) 과거 모델의 경험과 지식을 조직 자산으로 축적할 수 있다.

특히 생성형 AI 시대에는 LLM 파인튜닝 모델, RAG 시스템, 프롬프트 템플릿 등 새로운 유형의 아티팩트가 Retire 대상에 추가되었으며, 기반 모델 제공사의 Deprecation 일정에 따른 ``강제 Retire'' 시나리오도 대비해야 한다. EU AI Act의 10년 문서 보존 의무, GDPR의 삭제권과 보존 의무 간 긴장 관계 등 규제 환경의 복잡성이 증가하는 만큼, 본 장에서 제시한 프레임워크와 자동화 도구를 조직의 상황에 맞게 적용하여 AI 거버넌스의 완결성을 확보하기를 권고한다.
